\acknowledgement

回望整个项目的开发历程，从最初面对空白编辑器时的茫然，到逐步厘清系统架构、打通各个功能模块，再到论文撰写过程中对每一个技术细节的反复推敲，这段经历远比预想中更漫长，也更充实。

首先要感谢指导教师在选题和研究方向上给予的悉心指引。每一次讨论都帮助厘清了方向上的困惑，每一条修改意见都推动着系统设计走向更完善。没有导师的耐心指导，这个系统不会是现在的样子。

项目开发过程中遇到的最大挑战，莫过于 AI 流式对话功能的调试。SSE 的长连接管理、DeepSeek API 响应流的分块解析、中文字符在多字节边界处的截断处理——每一个问题都花费了相当多的时间排查。正是在一次次报错和修复的循环中，对前后端协作机制和 Node.js 异步模型的理解逐渐从理论走向了实践。

感谢同学们在开发过程中提供的反馈与建议。他们作为第一批使用者，指出了若干操作流程上的不便之处，推动了界面交互的多次迭代改进。

感谢家人在整个大学阶段给予的无条件支持与包容。无论是深夜调试代码时的陪伴，还是论文写作进展迟缓时的鼓励，都是完成这项工作不可或缺的动力来源。

本系统仍有诸多不足之处，如密码找回、实时推送、移动端适配等功能尚待完善。这些不足既是遗憾，也是继续探索的起点。希望这份工作能为后续从事相关系统开发的同学提供一些参考，也期待自己在未来的工作中将这些想法一一实现。

